\documentclass[10pt,a4paper]{amsart}
\usepackage{amsmath,amssymb,amsthm,graphicx}
\usepackage{verbatim}
\newenvironment{reference}{\comment}{\endcomment}
%\newenvironment{reference}{}{}
\newenvironment{slogan}{\comment}{\endcomment}
\newenvironment{history}{\comment}{\endcomment}

% For commutative diagrams we use Xy-pic
\usepackage[all]{xy}

% We use 2cell for 2-commutative diagrams.
\xyoption{2cell}
% Theorem environments.
%
\theoremstyle{plain}
\newtheorem{theorem}[subsection]{Theorem}
\newtheorem{proposition}[subsection]{Proposition}
\newtheorem{lemma}[subsection]{Lemma}

\theoremstyle{definition}
\newtheorem{definition}[subsection]{Definition}
\newtheorem{example}[subsection]{Example}
\newtheorem{exercise}[subsection]{Exercise}
\newtheorem{situation}[subsection]{Situation}

\theoremstyle{remark}
\newtheorem{remark}[subsection]{Remark}
\newtheorem{remarks}[subsection]{Remarks}

\numberwithin{equation}{subsection}

% Macros
%
\def\lim{\mathop{\mathrm{lim}}\nolimits}
\def\colim{\mathop{\mathrm{colim}}\nolimits}
\def\Spec{\mathop{\mathrm{Spec}}}
\def\Hom{\mathop{\mathrm{Hom}}\nolimits}
\def\Ext{\mathop{\mathrm{Ext}}\nolimits}
\def\SheafHom{\mathop{\mathcal{H}\!\mathit{om}}\nolimits}
\def\SheafExt{\mathop{\mathcal{E}\!\mathit{xt}}\nolimits}
\def\Sch{\mathit{Sch}}
\def\Mor{\mathop{\mathrm{Mor}}\nolimits}
\def\Ob{\mathop{\mathrm{Ob}}\nolimits}
\def\Sh{\mathop{\mathit{Sh}}\nolimits}
\def\NL{\mathop{N\!L}\nolimits}
\def\CH{\mathop{\mathrm{CH}}\nolimits}
\def\proetale{{pro\text{-}\acute{e}tale}}
\def\etale{{\acute{e}tale}}
\def\QCoh{\mathit{QCoh}}
\def\Ker{\mathop{\mathrm{Ker}}}
\def\Im{\mathop{\mathrm{Im}}}
\def\Coker{\mathop{\mathrm{Coker}}}
\def\Coim{\mathop{\mathrm{Coim}}}

% Boxtimes
%
\DeclareMathSymbol{\boxtimes}{\mathbin}{AMSa}{"02}

%
% Macros for moduli stacks/spaces
%
\def\QCohstack{\mathcal{QC}\!\mathit{oh}}
\def\Cohstack{\mathcal{C}\!\mathit{oh}}
\def\Spacesstack{\mathcal{S}\!\mathit{paces}}
\def\Quotfunctor{\mathrm{Quot}}
\def\Hilbfunctor{\mathrm{Hilb}}
\def\Curvesstack{\mathcal{C}\!\mathit{urves}}
\def\Polarizedstack{\mathcal{P}\!\mathit{olarized}}
\def\Complexesstack{\mathcal{C}\!\mathit{omplexes}}
% \Pic is the operator that assigns to X its picard group, usage \Pic(X)
% \Picardstack_{X/B} denotes the Picard stack of X over B
% \Picardfunctor_{X/B} denotes the Picard functor of X over B
\def\Pic{\mathop{\mathrm{Pic}}\nolimits}
\def\Picardstack{\mathcal{P}\!\mathit{ic}}
\def\Picardfunctor{\mathrm{Pic}}
\def\Deformationcategory{\mathcal{D}\!\mathit{ef}}
\usepackage{cleveref}
\crefname{originalreference}{original reference}{original references}
\Crefname{originalreference}{Original reference}{Original references}
\makeatletter
\newlabel{lemma-diagonal-projective-bundle}{{0FGR}{source}{}{sourcepointer-1}{}}
\newlabel{lemma-diagonal-projective-bundle@cref}{{[originalreference][0][]0FGR}{[1][0][]source}{}{}{}}
\makeatother
\crefname{originalreference}{original reference}{original references}
\Crefname{originalreference}{Original reference}{Original references}
\makeatletter
\newlabel{lemma-pushforward}{{03M9}{source}{}{sourcepointer-2}{}}
\newlabel{lemma-pushforward@cref}{{[originalreference][0][]03M9}{[1][0][]source}{}{}{}}
\makeatother
\title{The cycle class map commutes with pushforward along a projective bundle}
\author{The Stacks Project Authors}
\date{stacks-project:weil.tex:local-snapshot-20261012; GNU FDL 1.2 or later}
\begin{document}
\maketitle
\noindent\textit{Editorial note.} All delivered mathematical text is the author's own, copied from the pinned Stacks Project chapter file weil.tex; no statement or proof step was rewritten, reordered or removed.\par\medskip
\noindent\textit{Editorial note.} The counted proof is the author's complete argument: the reduction to the standard generators of the Chow groups of the projective bundle, the projection formula step, and the verification of the degree map through the class of the diagonal.\par\medskip
\noindent\textit{Editorial note.} The label lemma-ok-for-projective-bundle names the counted statement, and the proof environment immediately following it in the source is delivered verbatim as its proof.\par\medskip
\noindent\textit{Editorial note.} The counted statement is carried by the author's own lemma environment in the chapter source.\par\medskip
\noindent\textit{Editorial note.} Pre-open gate: python operations/agent-authoring/tools/precard\_gate.py --source weil.tex --label lemma-ok-for-projective-bundle --start 5025 --end 5061 --local-id 28146 returned RESULT=OK (L28146 free, occupancy FREE, public interval clear). Pre-assembly tag self-check: canonical tag resolved and not present in reference/public-items.jsonl.\par\medskip
\noindent\textit{Editorial note.} The statement and proof are copied from the pinned chapter file; the cited inputs are kept as external pointers with their canonical Stacks tags.\par\medskip
\noindent\textit{Allowed standard prerequisites (editorial).}\begin{itemize}
\item The projective space bundle formula for the cohomology theory. Axiom (A3) describing the cohomology of the projectivization of a vector bundle together with the splitting used to describe the pushforward of the powers of the twist. Identifies the generators of the Chow groups of the projective bundle and the pushforward of the powers of the first Chern class of the twist.
\item The projection formula for the cohomology theory. Same chapter, the lemma stating the projection formula, and the compatibility of the cycle class map with pullbacks. Moves the pushforward inside the product with the pullback of the class of the base cycle.
\item The computation of the class of the diagonal of a projective bundle. Same chapter, the explicit description of the class of the diagonal of a projective bundle. Verifies the condition characterizing the degree map of the projective bundle, which is the remaining input.
\end{itemize}\medskip
\ifdefined\hypertarget\hypertarget{sourcepointer-1}{}\fi
\noindent\textit{Source reference.} Original reference 0FGR: sources/stacks-weil/pinned-original-source.tex, source lines 1466--1466. Prerequisite pointer invoked at proof line(s) 5060; it names a standard fact of the same chapter that the counted argument uses, and every step of the counted proof built on it is retained verbatim in the delivered body.\par\medskip
\ifdefined\hypertarget\hypertarget{sourcepointer-2}{}\fi
\noindent\textit{Source reference.} Original reference 03M9: sources/stacks-weil/pinned-original-source.tex, source lines 2660--2660. Prerequisite pointer invoked at proof line(s) 5040; it names a standard fact of the same chapter that the counted argument uses, and every step of the counted proof built on it is retained verbatim in the delivered body.\par\medskip
\begin{lemma}
\label{lemma-ok-for-projective-bundle}
Assume given data (D0), (D1), and (D2') satisfying axioms (A1) -- (A7).
Let $p : P \to X$ be as in axiom (A3) with $X$ nonempty equidimensional.
Then $\gamma$ commutes with pushforward along $p$.
\end{lemma}

\begin{proof}
It suffices to prove this on generators for $\CH_*(P)$.
Thus it suffices to prove this for a cycle class of the
form $\xi^i \cdot p^*\alpha$ where $0 \leq i \leq r - 1$
and $\alpha \in \CH_*(X)$. Note that $p_*(\xi^i \cdot p^*\alpha) = 0$
if $i < r - 1$ and $p_*(\xi^{r - 1} \cdot p^*\alpha) = \alpha$.
On the other hand, we have
$\gamma(\xi^i \cdot p^*\alpha) = c^i \cup p^*\gamma(\alpha)$
and by the projection formula (Lemma \ref{lemma-pushforward})
we have
$$
p_*\gamma(\xi^i \cdot p^*\alpha) = p_*(c^i) \cup \gamma(\alpha)
$$
Thus it suffices to show that $p_*c^i = 0$ for $i < r - 1$ and
$p_*c^{r - 1} = 1$. Equivalently, it suffices to prove that
$\lambda_P : H^{2d + 2r - 2}(P)(d + r - 1) \to F$ defined by
the rules
$$
\lambda_P(c^i \cup p^*(a)) =
\left\{
\begin{matrix}
0 & \text{if} & i < r - 1 \\
\int_X(a) & \text{if} & i = r - 1
\end{matrix}
\right.
$$
satisfies the condition of axiom (A5). This follows from the
computation of the class of the diagonal of $P$ in
Lemma \ref{lemma-diagonal-projective-bundle}.
\end{proof}

\end{document}
